% Options for packages loaded elsewhere
\PassOptionsToPackage{unicode}{hyperref}
\PassOptionsToPackage{hyphens}{url}
\PassOptionsToPackage{dvipsnames,svgnames,x11names}{xcolor}
\documentclass[
  ngerman,
  10pt,
  a4paper,
]{article}
\usepackage{xcolor}
\usepackage[margin=22mm]{geometry}
\usepackage{amsmath,amssymb}
\setcounter{secnumdepth}{-\maxdimen} % remove section numbering
\usepackage{iftex}
\ifPDFTeX
  \usepackage[T1]{fontenc}
  \usepackage[utf8]{inputenc}
  \usepackage{textcomp} % provide euro and other symbols
\else % if luatex or xetex
  \usepackage{unicode-math} % this also loads fontspec
  \defaultfontfeatures{Scale=MatchLowercase}
  \defaultfontfeatures[\rmfamily]{Ligatures=TeX,Scale=1}
\fi
\usepackage{lmodern}
\ifPDFTeX\else
  % xetex/luatex font selection
  \setmainfont[Path=/Users/stefanhamann/.cache/codex-runtimes/codex-primary-runtime/dependencies/native/libreoffice-headless/libreoffice/LibreOfficeDev.app/Contents/Resources/fonts/truetype/,Extension=.ttf,BoldFont=DejaVuSans-Bold,ItalicFont=DejaVuSans-Oblique,BoldItalicFont=DejaVuSans-BoldOblique]{DejaVuSans}
  \setmonofont[Path=/Users/stefanhamann/.cache/codex-runtimes/codex-primary-runtime/dependencies/native/libreoffice-headless/libreoffice/LibreOfficeDev.app/Contents/Resources/fonts/truetype/,Extension=.ttf,BoldFont=DejaVuSansMono-Bold,ItalicFont=DejaVuSansMono-Oblique,BoldItalicFont=DejaVuSansMono-BoldOblique]{DejaVuSansMono}
\fi
% Use upquote if available, for straight quotes in verbatim environments
\IfFileExists{upquote.sty}{\usepackage{upquote}}{}
\IfFileExists{microtype.sty}{% use microtype if available
  \usepackage[]{microtype}
  \UseMicrotypeSet[protrusion]{basicmath} % disable protrusion for tt fonts
}{}
\makeatletter
\@ifundefined{KOMAClassName}{% if non-KOMA class
  \IfFileExists{parskip.sty}{%
    \usepackage{parskip}
  }{% else
    \setlength{\parindent}{0pt}
    \setlength{\parskip}{6pt plus 2pt minus 1pt}}
}{% if KOMA class
  \KOMAoptions{parskip=half}}
\makeatother
\usepackage{longtable,booktabs,array}
\newcounter{none} % for unnumbered tables
\usepackage{calc} % for calculating minipage widths
% Correct order of tables after \paragraph or \subparagraph
\usepackage{etoolbox}
\makeatletter
\patchcmd\longtable{\par}{\if@noskipsec\mbox{}\fi\par}{}{}
\makeatother
% Allow footnotes in longtable head/foot
\IfFileExists{footnotehyper.sty}{\usepackage{footnotehyper}}{\usepackage{footnote}}
\makesavenoteenv{longtable}
\ifLuaTeX
\usepackage[bidi=basic,shorthands=off]{babel}
\else
\usepackage[bidi=default,shorthands=off]{babel}
\fi
\ifPDFTeX
\else
\babelfont{rm}[Path=/Users/stefanhamann/.cache/codex-runtimes/codex-primary-runtime/dependencies/native/libreoffice-headless/libreoffice/LibreOfficeDev.app/Contents/Resources/fonts/truetype/,Extension=.ttf,BoldFont=DejaVuSans-Bold,ItalicFont=DejaVuSans-Oblique,BoldItalicFont=DejaVuSans-BoldOblique]{DejaVuSans}
\fi
\ifLuaTeX
  \usepackage{selnolig} % disable illegal ligatures
\fi
\setlength{\emergencystretch}{3em} % prevent overfull lines
\providecommand{\tightlist}{%
  \setlength{\itemsep}{0pt}\setlength{\parskip}{0pt}}
\usepackage{fvextra}
\RecustomVerbatimEnvironment{verbatim}{Verbatim}{breaklines=true,fontsize=\footnotesize}
\usepackage{xurl}
\usepackage{seqsplit}
\usepackage{adjustbox}
\usepackage{ragged2e}
\usepackage{titlesec}
\titleformat{\section}{\large\bfseries\raggedright}{}{0em}{}
\titleformat{\subsection}{\normalsize\bfseries\raggedright}{}{0em}{}
\DeclareRobustCommand{\codeword}[1]{\texttt{\seqsplit{#1}}}
\AtBeginEnvironment{longtable}{\small\RaggedRight}
\usepackage{newunicodechar}
\newunicodechar{𝒞}{\ensuremath{\mathcal C}}
\newunicodechar{𝔊}{\ensuremath{\mathfrak G}}
\newunicodechar{𝔤}{\ensuremath{\mathfrak g}}
\newunicodechar{𝓗}{\ensuremath{\mathcal H}}
\newunicodechar{ℋ}{\ensuremath{\mathcal H}}
\newunicodechar{𝔄}{\ensuremath{\mathfrak A}}
\newunicodechar{𝟙}{\ensuremath{\mathbf 1}}
\newunicodechar{𝕀}{\ensuremath{I}}
\newunicodechar{𝕋}{\ensuremath{\mathbb T}}
\newunicodechar{𝟘}{\ensuremath{\mathbf 0}}
\newunicodechar{⊞}{\ensuremath{\boxplus}}
\newunicodechar{∘}{\ensuremath{\circ}}
\newunicodechar{⟦}{\ensuremath{[\![}}
\newunicodechar{⟧}{\ensuremath{]\!]}}
\setlength{\emergencystretch}{4em}
\setlength{\parindent}{0pt}
\setlength{\parskip}{0.45em}
\AtBeginDocument{\hypersetup{colorlinks=true,linkcolor=black,urlcolor=blue}\pdfstringdefDisableCommands{\def\codeword#1{#1}}\RaggedRight}
\usepackage{fancyhdr}
\pagestyle{fancy}
\fancyhf{}
\fancyhead[L]{\small TFPT / Universalraum - v1.6.9}
\fancyfoot[R]{\thepage}
\setlength{\headheight}{15pt}
% The preserved full history now has three-digit page numbers.
\makeatletter
\renewcommand{\@pnumwidth}{2.6em}
\renewcommand{\@tocrmarg}{3.6em}
\makeatother

\usepackage{bookmark}
\IfFileExists{xurl.sty}{\usepackage{xurl}}{} % add URL line breaks if available
\urlstyle{same}
\hypersetup{
  pdflang={de-DE},
  colorlinks=true,
  linkcolor={Maroon},
  filecolor={Maroon},
  citecolor={Blue},
  urlcolor={Blue},
  pdfcreator={LaTeX via pandoc}}

\author{}
\date{}

\begin{document}

\section{Die Verbindung hinter den Schatten - einfach
erklärt}\label{die-verbindung-hinter-den-schatten---einfach-erkluxe4rt}

TFPT / Universalraum · 15. September 2026 · v1.6.9

\subsection{Der wichtigste neue Punkt}\label{der-wichtigste-neue-punkt}

Wir haben nicht die ganze Weltformel gefunden. Aber wir können jetzt an
zwei kleinen, genau gerechneten Beispielen zeigen, \textbf{wie eine
verborgene Verbindung aus ihrer Reaktion sichtbar wird}. Das ist ein
konkreter Schritt in die Richtung deiner Schatten-Idee.

Bislang hatten wir einen gut beschriebenen inneren Baustein. Jetzt
können wir genauer sagen, welche Information noch fehlt, wenn man aus
mehreren solchen Bausteinen ein gemeinsames System machen will - und
welche Reaktion diese Information verrät.

\subsection{Stell dir eine kleine Kupplung
vor}\label{stell-dir-eine-kleine-kupplung-vor}

Der ursprüngliche Mechanismus kann zwei Fermionen in ein Boson umwandeln
und wieder zurück. Seine innere Struktur ist sehr festgelegt. Die vielen
Einträge bedeuten dabei nicht viele unabhängige Maschinen: Das
ursprüngliche Modell benutzt bereits \textbf{64 gemeinsame
Fermionmoden}.

Die nächste Frage ist so einfach wie wichtig: Greifen zwei solche
Kupplungen auf dieselben Teile zu? Vollständig, teilweise oder gar
nicht?

In unserem kleinsten Vergleich lässt sich dieses gemeinsame Benutzen
durch eine Zahl beschreiben: die Überlappung. Beide Kupplungen können
einzeln vollkommen gleich aussehen, obwohl ihr gemeinsamer Ablauf
verschieden ist.

Bildlich: Zwei Wasserhähne sehen gleich aus. Ob sie an getrennten
Leitungen oder an teilweise gemeinsamen Leitungen hängen, erkennt man
nicht am Metall des Hahns. Man muss prüfen, wie ein Eingriff am einen
die Antwort des anderen verändert. Das ist nur ein Bild; unsere Rechnung
benutzt Quantenamplituden, keine Wasserströmung.

\subsection{Was wir jetzt wirklich herausgerechnet
haben}\label{was-wir-jetzt-wirklich-herausgerechnet-haben}

\textbf{Erstens: Die Reaktion verrät die Verbindung.} Im
Zwei-Kupplungs-Modell stimmen die ersten drei Energiemomente überein.
Das vierte verrät die Überlappung eindeutig - innerhalb dieser
festgelegten Modellfamilie.

Eine zweite Rechnung ist noch deutlicher: Zwei verschieden verbundene
Modelle haben das ganze gleiche Zweiladungsspektrum. Erst der
Dreiladungsversuch unterscheidet sie. Aus der vollständig aufgelösten
höheren Antwort lässt sich die positive Kopplungsmatrix zurückgewinnen.

Das ist der neue brauchbare Kern der Schattenidee: \textbf{Nicht jedes
Bild enthält genug Information. Mehrere richtig gewählte Antworten
können die zusätzliche Struktur aufdecken.} Noch nicht bewiesen ist,
dass alle diese Antworten in unserer Realität oder direkt im
ursprünglichen Compiler verfügbar sind.

\textbf{Zweitens: Ein Handgriff wurde einfacher.} Für den kleinen
Eingriffsversuch benötigen wir nicht mehr den alten zusätzlichen Zugriff
auf eine einzelne Fermionmode. Eine Phase auf der Bosonzahl genügt. Erst
nach der weiteren Entwicklung verändert sich die Antwort einer
Fermionmode. Damit haben wir einen echten Eingriff und seine Folge
berechnet, nicht nur eine Korrelation.

\textbf{Drittens: Die Aufzeichnung ist Teil des Ablaufs.} Ein kleiner
Zeiger merkt sich nur „Paar oder Boson``. Er verrät nicht, welchen der
interferierenden Paarwege das System genommen hat. Deshalb bleibt diese
wichtige innere Interferenz erhalten. Auch die Wirkung des Zeigers und
seine Energieänderung sind mitgerechnet.

Der Zeiger, sein Koppler und die nötige Steuerarbeit kommen aber noch
nicht kostenlos aus dem Compiler. Das ist die Stelle, an der ein
berechneter Laborversuch noch kein sich selbst ausführendes Universum
ist.

\textbf{Viertens: Ein verlockender Teilchenweg wurde früh geprüft.} Wir
haben eine passende mathematische Übersetzung in einen relativistischen
Spinor gefunden. Aber im naheliegenden freien masselosen Ansatz ist das
übersetzte Signal exakt null. Der Kanal sieht richtig aus, trägt dort
jedoch kein sichtbares Teilchen.

Das ist wichtig: Es verhindert, auf einer bloß passenden Form jahrelang
weiterzubauen. Als Nächstes muss derselbe korrekte Kanal unter einer
tatsächlich begründeten Wechselwirkung oder anderen Feldzuordnung eine
nichtverschwindende Antwort liefern.

\subsection{Was hat sich gegenüber v1.6.8
verändert?}\label{was-hat-sich-gegenuxfcber-v1.6.8-veruxe4ndert}

{\def\LTcaptype{none} % do not increment counter
\begin{longtable}[]{@{}
  >{\raggedright\arraybackslash}p{(\linewidth - 2\tabcolsep) * \real{0.5000}}
  >{\raggedright\arraybackslash}p{(\linewidth - 2\tabcolsep) * \real{0.5000}}@{}}
\toprule\noalign{}
\begin{minipage}[b]{\linewidth}\raggedright
Vorher
\end{minipage} & \begin{minipage}[b]{\linewidth}\raggedright
Jetzt
\end{minipage} \\
\midrule\noalign{}
\endhead
\bottomrule\noalign{}
\endlastfoot
„Uns fehlt die gemeinsame Verbindung.`` & Zwei kleine vollständige
Vergleiche zeigen, was sie verändert und wie man sie zurückliest. \\
Zusätzlicher Einzelmoden-Handgriff im Versuch & Derselbe innere Prozess
funktioniert mit dem schon dokumentierten Bosonzahl-Kontrolltyp. \\
Aufzeichnung musste die Interferenzgrenze beachten & Ein konkreter
Zeiger und sein vollständiger Folgeeffekt samt Arbeitsbilanz sind
berechnet. \\
Gemischthändiger Feldkanal hatte eine Übersetzungslücke & Die minimale
Impulsübersetzung ist da; ihr freier masseloser Ansatz besteht den
Nichtnulltest nicht. \\
Clock als möglicher Weg zu mehreren Vorkommen & Der geprüfte
ursprüngliche Clock bleibt im selben hellen Raum; er kopiert ihn
nicht. \\
\end{longtable}
}

\subsection{Wie könnte das gesuchte Gesamtbild
aussehen?}\label{wie-kuxf6nnte-das-gesuchte-gesamtbild-aussehen}

Eine einzige ursprüngliche Regel legt fest, \textbf{was gemeinsam
benutzt wird und was miteinander wechselwirken kann}. Aus derselben
Regel und demselben Zustand entstehen dann verschiedene Ansichten:

\begin{itemize}
\tightlist
\item
  Zeit wäre durch konkrete Entwicklungen und Aufzeichnungen beschrieben.
\item
  Geometrie müsste aus der Struktur tatsächlich unterscheidbarer und
  beeinflussbarer Teile folgen.
\item
  Teilchen müssten stabile, richtig transformierende Anregungen
  derselben Dynamik sein.
\item
  Konstanten müssten aus dieser Quelle berechnet und eindeutig auf
  Messgrößen abgebildet werden.
\end{itemize}

Das ist die gesuchte gemeinsame Bauidee, \textbf{noch nicht das
bewiesene Ergebnis}. Insbesondere sind Raumdimension, Relativität,
chirale Materie und Gravitation dadurch nicht automatisch erledigt. Auch
eine holographische Deutung benötigt eine eigene genaue Abbildung, nicht
nur ähnliche Bilder.

\subsection{Der kürzeste sinnvolle nächste
Weg}\label{der-kuxfcrzeste-sinnvolle-nuxe4chste-weg}

\textbf{1. Zwei ursprüngliche Teile und ihre Verbindung festnageln.} Aus
den bereits vorhandenen Quellobjekten muss folgen, welche zwei Teile
gemeinsam Ressourcen benutzen und wie stark. Nicht die schöne Zahl
auswählen, sondern die vollständige Zuordnung herleiten. Unsere neuen
Antwortformeln prüfen sofort, ob sie etwas Bestimmtes vorhersagt.

\textbf{2. Den Handgriff und den Zeiger in dieselbe Maschine einbauen.}
Dann müssen wir nicht von außen den passenden Puls setzen. Der
gemeinsame Anfangszustand und die gemeinsame Dynamik müssen Puls,
Aufzeichnung, Arbeit und Folgeantwort erzeugen.

\textbf{3. Ein wirklich sichtbares Teilchen auf genau dieser Grundlage
zeigen.} Zuerst muss der korrekte Feldkanal nichtnull sein. Danach lohnt
sich die Frage, ob viele solche Teile die richtige Raumzeit und Physik
ergeben.

Diese Runde hat die vollständige TOE nicht bewiesen. Sie hat aber aus
„irgendetwas fehlt`` eine kleinere, testbare Frage gemacht:
\textbf{Welche ursprüngliche Regel bestimmt die gemeinsame Benutzung -
und stimmt ihre vorhergesagte Antwort mit allen anderen Ansichten
überein?}

Die Rechnungen bestanden 514 exakt arbeitende Prüfbedingungen und vier
numerische Gegenkontrollen pro Ausführungsart. Das belegt die
angegebenen kleinen Modelle, nicht unser gesamtes Universum.

\end{document}
